\documentclass[a4paper,11pt]{article}
\usepackage[utf8]{inputenc}
\usepackage[T1]{fontenc}
\usepackage[french]{babel}
\usepackage{lmodern}
\usepackage{amsmath,amssymb,amsthm}
\usepackage{mathrsfs}
\usepackage{enumitem}
\usepackage{multicol}

\setlist[enumerate]{label=\textbf{\textcolor{Couleur8}{\arabic*.}}, left=1em}

% Si vous voulez aussi modifier les listes enumerate imbriquées :
\setlist[enumerate,2]{label=\textcolor{Couleur8}{\alph*)}}
\setlist[enumerate,3]{label=\textcolor{Couleur8}{\roman*)}}
\setlist[enumerate,4]{label=\textcolor{Couleur8}{\Alph*)}}
\usepackage{graphicx}
\usepackage{xcolor}
\usepackage{tcolorbox}
\usepackage{geometry}
\usepackage{fancyhdr}
\usepackage{titlesec}
\usepackage{cancel}
\usepackage{ifthen}
\usepackage{tikz}
\usetikzlibrary{arrows.meta}
\usepackage{tkz-tab}
\usepackage{eurosym}

% OPTION PROF/ÉLÈVE
% Décommentez UNE des deux lignes suivantes :
\newboolean{prof}\setboolean{prof}{true}   % Version professeur (avec corrections des devoirs)
\newboolean{prof}\setboolean{prof}{true}  % Version élève (sans corrections des devoirs)

\definecolor{Couleur1}{HTML}{4A5D7A} %Th
\definecolor{Couleur2}{HTML}{A5B5D6} %Prop
\definecolor{Couleur3}{HTML}{A8C68A} %Méthode
\definecolor{Couleur6}{HTML}{8A9CC1} %Def
\definecolor{Couleur7}{HTML}{E6B459} %Rq Voc
\definecolor{Couleur8}{HTML}{D36740} %Titres
\definecolor{correction}{HTML}{D36740} % Couleur pour les corrections
\definecolor{coopmathscorr}{HTML}{F15929} % Couleur corrections coopmaths

% Configuration des marges
\geometry{
	top=2cm,
	bottom=2cm,
	inner=1.5cm,
	outer=1.5cm,
}

% Police
\renewcommand{\familydefault}{\sfdefault}

% Compteurs
\newcounter{seancenum}
\newcounter{seancenumD}
\setcounter{seancenum}{1}
\setcounter{seancenumD}{1}

% Configuration des en-têtes et pieds de page
\pagestyle{fancy}
\fancyhf{}
\ifthenelse{\boolean{prof}}{
    \fancyhead[L]{\textcolor{Couleur8}{\textbf{Corrections Automatismes - Classe de seconde - Version Professeur}}}
}{
    \fancyhead[L]{\textcolor{Couleur8}{\textbf{Corrections Devoirs - Séance 10 - Classe de seconde}}}
}
\fancyhead[R]{\textcolor{Couleur8}{\textbf{2025-2026}}}
\fancyfoot[L]{\textcolor{Couleur8}{Mme Bertail et Mme Pinto}}
\fancyfoot[R]{\textcolor{Couleur8}{\thepage}}
\renewcommand{\headrulewidth}{1pt}
\renewcommand{\headrule}{\hbox to\headwidth{\color{Couleur8}\leaders\hrule height \headrulewidth\hfill}}

% Environnements personnalisés
\newtcolorbox{seance}[1]{
    colback=Couleur6!10,
    colframe=Couleur1!65,
    fonttitle=\bfseries\large,
    title=Correction Séance \theseancenum #1,
    left=5mm,
    right=5mm,
    top=3mm,
    bottom=3mm,
    arc=0mm,
    after=\stepcounter{seancenum}
}

\newtcolorbox{seanceD}[1]{
    colback=Couleur6!5,
    colframe=Couleur6!45,
    fonttitle=\bfseries\large,
    title=Correction Devoirs - Séance \theseancenumD #1,
    left=5mm,
    right=5mm,
    top=3mm,
    bottom=3mm,
    arc=0mm,
    after=\stepcounter{seancenumD}
}

\begin{document}

\setcounter{seancenumD}{10}
\newpage
\begin{seanceD}{} %10
\begin{enumerate}
\item On remplace $x$ par $-\dfrac{2}{3}$ dans l'expression de $f$ :\\
     
    $\begin{aligned}
    f\left(-\dfrac{2}{3}\right)&=-2\times \left(-\dfrac{2}{3}\right)^2-2\times \left(-\dfrac{2}{3}\right)+1\\
    &=-2\times\dfrac{4}{9}+\dfrac{4}{3}+1\\
    &=-\dfrac{8}{9}+\dfrac{4}{3}+1\\
    &=\dfrac{-8+12+9}{9}\\
    &=\dfrac{13}{9}
    \end{aligned}$\\
    
    L'image de $-\dfrac{2}{3}$ par la fonction $f$ est : ${\color[HTML]{f15929}\boldsymbol{\dfrac{13}{9}}}$.

\item En prenant deux points sur la droite, on obtient :\\
     
    $m=\dfrac{{\color{blue}\boldsymbol{4}}}{{\color{red}\boldsymbol{3}}}={\color[HTML]{f15929}\boldsymbol{\dfrac{4}{3}}}$\\\begin{tikzpicture}[baseline,scale = 0.75]

    \tikzset{
      point/.style={
        thick,
        draw,
        cross out,
        inner sep=0pt,
        minimum width=5pt,
        minimum height=5pt,
      },
    }
    \clip (-6,-1) rectangle (6,6.25);
    	\draw[color={blue},line width = 2] (-33,-40)--(30,44);
	\draw[color ={black},line width = 1.2,->] (-6,0)--(6,0);
	\draw[color ={black},line width = 1.2,->] (0,-1)--(0,6);
	\draw[color ={black},opacity = 0.4] (-6,1)--(6,1);
	\draw[color ={black},opacity = 0.4] (-6,-1)--(6,-1);
	\draw[color ={black},opacity = 0.4] (-6,2)--(6,2);
	\draw[color ={black},opacity = 0.4] (-6,3)--(6,3);
	\draw[color ={black},opacity = 0.4] (-6,4)--(6,4);
	\draw[color ={black},opacity = 0.4] (-6,5)--(6,5);
	\draw[color ={black},opacity = 0.4] (1,-1)--(1,6);
	\draw[color ={black},opacity = 0.4] (-1,-1)--(-1,6);
	\draw[color ={black},opacity = 0.4] (2,-1)--(2,6);
	\draw[color ={black},opacity = 0.4] (-2,-1)--(-2,6);
	\draw[color ={black},opacity = 0.4] (3,-1)--(3,6);
	\draw[color ={black},opacity = 0.4] (-3,-1)--(-3,6);
	\draw[color ={black},opacity = 0.4] (4,-1)--(4,6);
	\draw[color ={black},opacity = 0.4] (-4,-1)--(-4,6);
	\draw[color ={black},opacity = 0.4] (5,-1)--(5,6);
	\draw[color ={black},opacity = 0.4] (-5,-1)--(-5,6);
	\draw[color ={gray},opacity = 0.1] (-6.1,0)--(6.1,0);
	\draw[color ={gray},opacity = 0.1] (-6.1,0.1)--(6.1,0.1);
	\draw[color ={gray},opacity = 0.1] (-6.1,-0.1)--(6.1,-0.1);
	\draw[color ={gray},opacity = 0.1] (-6.1,0.2)--(6.1,0.2);
	\draw[color ={gray},opacity = 0.1] (-6.1,-0.2)--(6.1,-0.2);
	\draw[color ={gray},opacity = 0.1] (-6.1,0.3)--(6.1,0.3);
	\draw[color ={gray},opacity = 0.1] (-6.1,-0.3)--(6.1,-0.3);
	\draw[color ={gray},opacity = 0.1] (-6.1,0.4)--(6.1,0.4);
	\draw[color ={gray},opacity = 0.1] (-6.1,-0.4)--(6.1,-0.4);
	\draw[color ={gray},opacity = 0.1] (-6.1,0.5)--(6.1,0.5);
	\draw[color ={gray},opacity = 0.1] (-6.1,-0.5)--(6.1,-0.5);
	\draw[color ={gray},opacity = 0.1] (-6.1,0.6)--(6.1,0.6);
	\draw[color ={gray},opacity = 0.1] (-6.1,-0.6)--(6.1,-0.6);
	\draw[color ={gray},opacity = 0.1] (-6.1,0.7)--(6.1,0.7);
	\draw[color ={gray},opacity = 0.1] (-6.1,-0.7)--(6.1,-0.7);
	\draw[color ={gray},opacity = 0.1] (-6.1,0.8)--(6.1,0.8);
	\draw[color ={gray},opacity = 0.1] (-6.1,-0.8)--(6.1,-0.8);
	\draw[color ={gray},opacity = 0.1] (-6.1,0.9)--(6.1,0.9);
	\draw[color ={gray},opacity = 0.1] (-6.1,-0.9)--(6.1,-0.9);
	\draw[color ={gray},opacity = 0.1] (-6.1,1)--(6.1,1);
	\draw[color ={gray},opacity = 0.1] (-6.1,-1)--(6.1,-1);
	\draw[color ={gray},opacity = 0.1] (-6.1,1.1)--(6.1,1.1);
	\draw[color ={gray},opacity = 0.1] (-6.1,-1.1)--(6.1,-1.1);
	\draw[color ={gray},opacity = 0.1] (-6.1,1.2)--(6.1,1.2);
	\draw[color ={gray},opacity = 0.1] (-6.1,1.3)--(6.1,1.3);
	\draw[color ={gray},opacity = 0.1] (-6.1,1.4)--(6.1,1.4);
	\draw[color ={gray},opacity = 0.1] (-6.1,1.5)--(6.1,1.5);
	\draw[color ={gray},opacity = 0.1] (-6.1,1.6)--(6.1,1.6);
	\draw[color ={gray},opacity = 0.1] (-6.1,1.7)--(6.1,1.7);
	\draw[color ={gray},opacity = 0.1] (-6.1,1.8)--(6.1,1.8);
	\draw[color ={gray},opacity = 0.1] (-6.1,1.9)--(6.1,1.9);
	\draw[color ={gray},opacity = 0.1] (-6.1,2)--(6.1,2);
	\draw[color ={gray},opacity = 0.1] (-6.1,2.1)--(6.1,2.1);
	\draw[color ={gray},opacity = 0.1] (-6.1,2.2)--(6.1,2.2);
	\draw[color ={gray},opacity = 0.1] (-6.1,2.3)--(6.1,2.3);
	\draw[color ={gray},opacity = 0.1] (-6.1,2.4)--(6.1,2.4);
	\draw[color ={gray},opacity = 0.1] (-6.1,2.5)--(6.1,2.5);
	\draw[color ={gray},opacity = 0.1] (-6.1,2.6)--(6.1,2.6);
	\draw[color ={gray},opacity = 0.1] (-6.1,2.7)--(6.1,2.7);
	\draw[color ={gray},opacity = 0.1] (-6.1,2.8)--(6.1,2.8);
	\draw[color ={gray},opacity = 0.1] (-6.1,2.9)--(6.1,2.9);
	\draw[color ={gray},opacity = 0.1] (-6.1,3)--(6.1,3);
	\draw[color ={gray},opacity = 0.1] (-6.1,3.1)--(6.1,3.1);
	\draw[color ={gray},opacity = 0.1] (-6.1,3.2)--(6.1,3.2);
	\draw[color ={gray},opacity = 0.1] (-6.1,3.3)--(6.1,3.3);
	\draw[color ={gray},opacity = 0.1] (-6.1,3.4)--(6.1,3.4);
	\draw[color ={gray},opacity = 0.1] (-6.1,3.5)--(6.1,3.5);
	\draw[color ={gray},opacity = 0.1] (-6.1,3.6)--(6.1,3.6);
	\draw[color ={gray},opacity = 0.1] (-6.1,3.7)--(6.1,3.7);
	\draw[color ={gray},opacity = 0.1] (-6.1,3.8)--(6.1,3.8);
	\draw[color ={gray},opacity = 0.1] (-6.1,3.9)--(6.1,3.9);
	\draw[color ={gray},opacity = 0.1] (-6.1,4)--(6.1,4);
	\draw[color ={gray},opacity = 0.1] (-6.1,4.1)--(6.1,4.1);
	\draw[color ={gray},opacity = 0.1] (-6.1,4.2)--(6.1,4.2);
	\draw[color ={gray},opacity = 0.1] (-6.1,4.3)--(6.1,4.3);
	\draw[color ={gray},opacity = 0.1] (-6.1,4.4)--(6.1,4.4);
	\draw[color ={gray},opacity = 0.1] (-6.1,4.5)--(6.1,4.5);
	\draw[color ={gray},opacity = 0.1] (-6.1,4.6)--(6.1,4.6);
	\draw[color ={gray},opacity = 0.1] (-6.1,4.7)--(6.1,4.7);
	\draw[color ={gray},opacity = 0.1] (-6.1,4.8)--(6.1,4.8);
	\draw[color ={gray},opacity = 0.1] (-6.1,4.9)--(6.1,4.9);
	\draw[color ={gray},opacity = 0.1] (-6.1,5)--(6.1,5);
	\draw[color ={gray},opacity = 0.1] (-6.1,5.1)--(6.1,5.1);
	\draw[color ={gray},opacity = 0.1] (-6.1,5.2)--(6.1,5.2);
	\draw[color ={gray},opacity = 0.1] (-6.1,5.3)--(6.1,5.3);
	\draw[color ={gray},opacity = 0.1] (-6.1,5.4)--(6.1,5.4);
	\draw[color ={gray},opacity = 0.1] (-6.1,5.5)--(6.1,5.5);
	\draw[color ={gray},opacity = 0.1] (-6.1,5.6)--(6.1,5.6);
	\draw[color ={gray},opacity = 0.1] (-6.1,5.7)--(6.1,5.7);
	\draw[color ={gray},opacity = 0.1] (-6.1,5.8)--(6.1,5.8);
	\draw[color ={gray},opacity = 0.1] (-6.1,5.9)--(6.1,5.9);
	\draw[color ={gray},opacity = 0.1] (-6.1,6)--(6.1,6);
	\draw[color ={gray},opacity = 0.1] (-6.1,6.1)--(6.1,6.1);
	\draw[color ={gray},opacity = 0.1] (0,-1.1)--(0,6.1);
	\draw[color ={gray},opacity = 0.1] (0.1,-1.1)--(0.1,6.1);
	\draw[color ={gray},opacity = 0.1] (-0.1,-1.1)--(-0.1,6.1);
	\draw[color ={gray},opacity = 0.1] (0.2,-1.1)--(0.2,6.1);
	\draw[color ={gray},opacity = 0.1] (-0.2,-1.1)--(-0.2,6.1);
	\draw[color ={gray},opacity = 0.1] (0.3,-1.1)--(0.3,6.1);
	\draw[color ={gray},opacity = 0.1] (-0.3,-1.1)--(-0.3,6.1);
	\draw[color ={gray},opacity = 0.1] (0.4,-1.1)--(0.4,6.1);
	\draw[color ={gray},opacity = 0.1] (-0.4,-1.1)--(-0.4,6.1);
	\draw[color ={gray},opacity = 0.1] (0.5,-1.1)--(0.5,6.1);
	\draw[color ={gray},opacity = 0.1] (-0.5,-1.1)--(-0.5,6.1);
	\draw[color ={gray},opacity = 0.1] (0.6,-1.1)--(0.6,6.1);
	\draw[color ={gray},opacity = 0.1] (-0.6,-1.1)--(-0.6,6.1);
	\draw[color ={gray},opacity = 0.1] (0.7,-1.1)--(0.7,6.1);
	\draw[color ={gray},opacity = 0.1] (-0.7,-1.1)--(-0.7,6.1);
	\draw[color ={gray},opacity = 0.1] (0.8,-1.1)--(0.8,6.1);
	\draw[color ={gray},opacity = 0.1] (-0.8,-1.1)--(-0.8,6.1);
	\draw[color ={gray},opacity = 0.1] (0.9,-1.1)--(0.9,6.1);
	\draw[color ={gray},opacity = 0.1] (-0.9,-1.1)--(-0.9,6.1);
	\draw[color ={gray},opacity = 0.1] (1,-1.1)--(1,6.1);
	\draw[color ={gray},opacity = 0.1] (-1,-1.1)--(-1,6.1);
	\draw[color ={gray},opacity = 0.1] (1.1,-1.1)--(1.1,6.1);
	\draw[color ={gray},opacity = 0.1] (-1.1,-1.1)--(-1.1,6.1);
	\draw[color ={gray},opacity = 0.1] (1.2,-1.1)--(1.2,6.1);
	\draw[color ={gray},opacity = 0.1] (-1.2,-1.1)--(-1.2,6.1);
	\draw[color ={gray},opacity = 0.1] (1.3,-1.1)--(1.3,6.1);
	\draw[color ={gray},opacity = 0.1] (-1.3,-1.1)--(-1.3,6.1);
	\draw[color ={gray},opacity = 0.1] (1.4,-1.1)--(1.4,6.1);
	\draw[color ={gray},opacity = 0.1] (-1.4,-1.1)--(-1.4,6.1);
	\draw[color ={gray},opacity = 0.1] (1.5,-1.1)--(1.5,6.1);
	\draw[color ={gray},opacity = 0.1] (-1.5,-1.1)--(-1.5,6.1);
	\draw[color ={gray},opacity = 0.1] (1.6,-1.1)--(1.6,6.1);
	\draw[color ={gray},opacity = 0.1] (-1.6,-1.1)--(-1.6,6.1);
	\draw[color ={gray},opacity = 0.1] (1.7,-1.1)--(1.7,6.1);
	\draw[color ={gray},opacity = 0.1] (-1.7,-1.1)--(-1.7,6.1);
	\draw[color ={gray},opacity = 0.1] (1.8,-1.1)--(1.8,6.1);
	\draw[color ={gray},opacity = 0.1] (-1.8,-1.1)--(-1.8,6.1);
	\draw[color ={gray},opacity = 0.1] (1.9,-1.1)--(1.9,6.1);
	\draw[color ={gray},opacity = 0.1] (-1.9,-1.1)--(-1.9,6.1);
	\draw[color ={gray},opacity = 0.1] (2,-1.1)--(2,6.1);
	\draw[color ={gray},opacity = 0.1] (-2,-1.1)--(-2,6.1);
	\draw[color ={gray},opacity = 0.1] (2.1,-1.1)--(2.1,6.1);
	\draw[color ={gray},opacity = 0.1] (-2.1,-1.1)--(-2.1,6.1);
	\draw[color ={gray},opacity = 0.1] (2.2,-1.1)--(2.2,6.1);
	\draw[color ={gray},opacity = 0.1] (-2.2,-1.1)--(-2.2,6.1);
	\draw[color ={gray},opacity = 0.1] (2.3,-1.1)--(2.3,6.1);
	\draw[color ={gray},opacity = 0.1] (-2.3,-1.1)--(-2.3,6.1);
	\draw[color ={gray},opacity = 0.1] (2.4,-1.1)--(2.4,6.1);
	\draw[color ={gray},opacity = 0.1] (-2.4,-1.1)--(-2.4,6.1);
	\draw[color ={gray},opacity = 0.1] (2.5,-1.1)--(2.5,6.1);
	\draw[color ={gray},opacity = 0.1] (-2.5,-1.1)--(-2.5,6.1);
	\draw[color ={gray},opacity = 0.1] (2.6,-1.1)--(2.6,6.1);
	\draw[color ={gray},opacity = 0.1] (-2.6,-1.1)--(-2.6,6.1);
	\draw[color ={gray},opacity = 0.1] (2.7,-1.1)--(2.7,6.1);
	\draw[color ={gray},opacity = 0.1] (-2.7,-1.1)--(-2.7,6.1);
	\draw[color ={gray},opacity = 0.1] (2.8,-1.1)--(2.8,6.1);
	\draw[color ={gray},opacity = 0.1] (-2.8,-1.1)--(-2.8,6.1);
	\draw[color ={gray},opacity = 0.1] (2.9,-1.1)--(2.9,6.1);
	\draw[color ={gray},opacity = 0.1] (-2.9,-1.1)--(-2.9,6.1);
	\draw[color ={gray},opacity = 0.1] (3,-1.1)--(3,6.1);
	\draw[color ={gray},opacity = 0.1] (-3,-1.1)--(-3,6.1);
	\draw[color ={gray},opacity = 0.1] (3.1,-1.1)--(3.1,6.1);
	\draw[color ={gray},opacity = 0.1] (-3.1,-1.1)--(-3.1,6.1);
	\draw[color ={gray},opacity = 0.1] (3.2,-1.1)--(3.2,6.1);
	\draw[color ={gray},opacity = 0.1] (-3.2,-1.1)--(-3.2,6.1);
	\draw[color ={gray},opacity = 0.1] (3.3,-1.1)--(3.3,6.1);
	\draw[color ={gray},opacity = 0.1] (-3.3,-1.1)--(-3.3,6.1);
	\draw[color ={gray},opacity = 0.1] (3.4,-1.1)--(3.4,6.1);
	\draw[color ={gray},opacity = 0.1] (-3.4,-1.1)--(-3.4,6.1);
	\draw[color ={gray},opacity = 0.1] (3.5,-1.1)--(3.5,6.1);
	\draw[color ={gray},opacity = 0.1] (-3.5,-1.1)--(-3.5,6.1);
	\draw[color ={gray},opacity = 0.1] (3.6,-1.1)--(3.6,6.1);
	\draw[color ={gray},opacity = 0.1] (-3.6,-1.1)--(-3.6,6.1);
	\draw[color ={gray},opacity = 0.1] (3.7,-1.1)--(3.7,6.1);
	\draw[color ={gray},opacity = 0.1] (-3.7,-1.1)--(-3.7,6.1);
	\draw[color ={gray},opacity = 0.1] (3.8,-1.1)--(3.8,6.1);
	\draw[color ={gray},opacity = 0.1] (-3.8,-1.1)--(-3.8,6.1);
	\draw[color ={gray},opacity = 0.1] (3.9,-1.1)--(3.9,6.1);
	\draw[color ={gray},opacity = 0.1] (-3.9,-1.1)--(-3.9,6.1);
	\draw[color ={gray},opacity = 0.1] (4,-1.1)--(4,6.1);
	\draw[color ={gray},opacity = 0.1] (-4,-1.1)--(-4,6.1);
	\draw[color ={gray},opacity = 0.1] (4.1,-1.1)--(4.1,6.1);
	\draw[color ={gray},opacity = 0.1] (-4.1,-1.1)--(-4.1,6.1);
	\draw[color ={gray},opacity = 0.1] (4.2,-1.1)--(4.2,6.1);
	\draw[color ={gray},opacity = 0.1] (-4.2,-1.1)--(-4.2,6.1);
	\draw[color ={gray},opacity = 0.1] (4.3,-1.1)--(4.3,6.1);
	\draw[color ={gray},opacity = 0.1] (-4.3,-1.1)--(-4.3,6.1);
	\draw[color ={gray},opacity = 0.1] (4.4,-1.1)--(4.4,6.1);
	\draw[color ={gray},opacity = 0.1] (-4.4,-1.1)--(-4.4,6.1);
	\draw[color ={gray},opacity = 0.1] (4.5,-1.1)--(4.5,6.1);
	\draw[color ={gray},opacity = 0.1] (-4.5,-1.1)--(-4.5,6.1);
	\draw[color ={gray},opacity = 0.1] (4.6,-1.1)--(4.6,6.1);
	\draw[color ={gray},opacity = 0.1] (-4.6,-1.1)--(-4.6,6.1);
	\draw[color ={gray},opacity = 0.1] (4.7,-1.1)--(4.7,6.1);
	\draw[color ={gray},opacity = 0.1] (-4.7,-1.1)--(-4.7,6.1);
	\draw[color ={gray},opacity = 0.1] (4.8,-1.1)--(4.8,6.1);
	\draw[color ={gray},opacity = 0.1] (-4.8,-1.1)--(-4.8,6.1);
	\draw[color ={gray},opacity = 0.1] (4.9,-1.1)--(4.9,6.1);
	\draw[color ={gray},opacity = 0.1] (-4.9,-1.1)--(-4.9,6.1);
	\draw[color ={gray},opacity = 0.1] (5,-1.1)--(5,6.1);
	\draw[color ={gray},opacity = 0.1] (-5,-1.1)--(-5,6.1);
	\draw[color ={gray},opacity = 0.1] (5.1,-1.1)--(5.1,6.1);
	\draw[color ={gray},opacity = 0.1] (-5.1,-1.1)--(-5.1,6.1);
	\draw[color ={gray},opacity = 0.1] (5.2,-1.1)--(5.2,6.1);
	\draw[color ={gray},opacity = 0.1] (-5.2,-1.1)--(-5.2,6.1);
	\draw[color ={gray},opacity = 0.1] (5.3,-1.1)--(5.3,6.1);
	\draw[color ={gray},opacity = 0.1] (-5.3,-1.1)--(-5.3,6.1);
	\draw[color ={gray},opacity = 0.1] (5.4,-1.1)--(5.4,6.1);
	\draw[color ={gray},opacity = 0.1] (-5.4,-1.1)--(-5.4,6.1);
	\draw[color ={gray},opacity = 0.1] (5.5,-1.1)--(5.5,6.1);
	\draw[color ={gray},opacity = 0.1] (-5.5,-1.1)--(-5.5,6.1);
	\draw[color ={gray},opacity = 0.1] (5.6,-1.1)--(5.6,6.1);
	\draw[color ={gray},opacity = 0.1] (-5.6,-1.1)--(-5.6,6.1);
	\draw[color ={gray},opacity = 0.1] (5.7,-1.1)--(5.7,6.1);
	\draw[color ={gray},opacity = 0.1] (-5.7,-1.1)--(-5.7,6.1);
	\draw[color ={gray},opacity = 0.1] (5.8,-1.1)--(5.8,6.1);
	\draw[color ={gray},opacity = 0.1] (-5.8,-1.1)--(-5.8,6.1);
	\draw[color ={gray},opacity = 0.1] (5.9,-1.1)--(5.9,6.1);
	\draw[color ={gray},opacity = 0.1] (-5.9,-1.1)--(-5.9,6.1);
	\draw[color ={gray},opacity = 0.1] (6,-1.1)--(6,6.1);
	\draw[color ={gray},opacity = 0.1] (-6,-1.1)--(-6,6.1);
	\draw[color ={gray},opacity = 0.1] (6.1,-1.1)--(6.1,6.1);
	\draw[color ={gray},opacity = 0.1] (-6.1,-1.1)--(-6.1,6.1);
	\draw[color ={black},line width = 1.2] (0,-0.1)--(0,0.1);
	\draw[color ={black},line width = 1.2] (1,-0.1)--(1,0.1);
	\draw[color ={black},line width = 1.2] (-1,-0.1)--(-1,0.1);
	\draw[color ={black},line width = 1.2] (2,-0.1)--(2,0.1);
	\draw[color ={black},line width = 1.2] (-2,-0.1)--(-2,0.1);
	\draw[color ={black},line width = 1.2] (3,-0.1)--(3,0.1);
	\draw[color ={black},line width = 1.2] (-3,-0.1)--(-3,0.1);
	\draw[color ={black},line width = 1.2] (4,-0.1)--(4,0.1);
	\draw[color ={black},line width = 1.2] (-4,-0.1)--(-4,0.1);
	\draw[color ={black},line width = 1.2] (-0.1,0)--(0.1,0);
	\draw[color ={black},line width = 1.2] (-0.1,1)--(0.1,1);
	\draw[color ={black},line width = 1.2] (-0.1,2)--(0.1,2);
	\draw[color ={black},line width = 1.2] (-0.1,3)--(0.1,3);
	\draw[color ={black},line width = 1.2] (-0.1,4)--(0.1,4);
	\draw[color ={black},line width = 1.2] (-0.1,5)--(0.1,5);
	\draw (1,-0.4) node[anchor = center, rotate=0] {\scriptsize \color{black}{$1$}};
	\draw (2,-0.4) node[anchor = center, rotate=0] {\scriptsize \color{black}{$2$}};
	\draw (3,-0.4) node[anchor = center, rotate=0] {\scriptsize \color{black}{$3$}};
	\draw (4,-0.4) node[anchor = center, rotate=0] {\scriptsize \color{black}{$4$}};
	\draw (5,-0.4) node[anchor = center, rotate=0] {\scriptsize \color{black}{$5$}};
	\draw (-1,-0.4) node[anchor = center, rotate=0] {\scriptsize \color{black}{$-1$}};
	\draw (-2,-0.4) node[anchor = center, rotate=0] {\scriptsize \color{black}{$-2$}};
	\draw (-3,-0.4) node[anchor = center, rotate=0] {\scriptsize \color{black}{$-3$}};
	\draw (-4,-0.4) node[anchor = center, rotate=0] {\scriptsize \color{black}{$-4$}};
	\draw (-5,-0.4) node[anchor = center, rotate=0] {\scriptsize \color{black}{$-5$}};
	\draw (-0.3,1.1) node[anchor = center, rotate=0] {\scriptsize \color{black}{$1$}};
	\draw (-0.3,2.1) node[anchor = center, rotate=0] {\scriptsize \color{black}{$2$}};
	\draw (-0.3,3.1) node[anchor = center, rotate=0] {\scriptsize \color{black}{$3$}};
	\draw (-0.3,4.1) node[anchor = center, rotate=0] {\scriptsize \color{black}{$4$}};
	\draw (-0.3,5.1) node[anchor = center, rotate=0] {\scriptsize \color{black}{$5$}};
	\draw[color ={black},line width = 1.25,opacity = 0.8] (-0.19,4.19)--(0.19,3.81);\draw[color ={black},line width = 1.25,opacity = 0.8] (-0.19,3.81)--(0.19,4.19);
	\draw (0,4.5) node[anchor = center, rotate=0] {\scriptsize \color{black}{$B$}};
	\draw (-3,0.5) node[anchor = center, rotate=0] {\scriptsize \color{black}{$A$}};
	\draw[color ={black},line width = 1.25,opacity = 0.8] (-3.19,0.19)--(-2.81,-0.19);\draw[color ={black},line width = 1.25,opacity = 0.8] (-3.19,-0.19)--(-2.81,0.19);
	\draw (-0.3,-0.3) node[anchor = center, rotate=0] {\scriptsize \color{black}{$\text{O}$}};
	\draw[color ={black},line width = 2, dashed ] (0,4)--(-3,4);
	\draw[color ={black},line width = 2, dashed ] (-3,0)--(-3,4);
	\draw (-1.5,4.3) node[anchor = center, rotate=0] {\scriptsize \color{red}{$-3$}};
	\draw (-2.5,2) node[anchor = center, rotate=0] {\scriptsize \color{blue}{$-4$}};

\end{tikzpicture}


\item  Multiplier par  $10^{-4}$ revient à multiplier par $0{,}000\,1$,  donc l'écriture décimale de $2{,}35\times 10^{-4}$ est : ${\color[HTML]{f15929}\boldsymbol{0{,}000\,235}}$.


\item  $A          =-4 + {\color{blue}\boldsymbol{3^3}} \times 3$\\$A=-4 + {\color{blue}\boldsymbol{27\times3}}$\\$A=-4+81$\\$A = {\color[HTML]{f15929}\boldsymbol{77}}$


\item 
 $-4x-6\geqslant -3x-3$\\On ajoute $3x$ aux deux membres.\\$-4x-6{\color{blue}\boldsymbol{+3x}}
          \geqslant-3x-3{\color{blue}\boldsymbol{+3x}}$\\$-x-6\geqslant-3$\\On ajoute $6$ aux deux membres.\\$-x-6{\color{blue}\boldsymbol{+6}}
          \geqslant-3{\color{blue}\boldsymbol{+6}}$\\$-x\geqslant3$\\On divise les deux membres par $-1$.\\Comme $-1$ est négatif, l'inégalité change de sens.\\$-x{\color{blue}\boldsymbol{\div(-1)\leqslant}}3{\color{blue}\boldsymbol{\div(-1)}}$\\$x\leqslant\dfrac{3}{-1}$\\$x\leqslant-3$\\ L'ensemble de solutions de l'inéquation est $S = {\color[HTML]{f15929}\boldsymbol{\left] -\infty\,;\,-3\right]}}$.

\end{enumerate}
\end{seanceD}

\end{document}